\section*{अध्याय \ref{ch:g7:digestion-and-nutrients} --- पाचन और पोषक तत्व}

\begin{solution}{exo:g7:digestion-and-nutrients:1}
वे भोजन की \emph{रासायनिक} तोड़-फोड़ करते हैं: उसके बड़े घटकों ---
स्टार्च, शरीर बनाने वाला पदार्थ, चिकनाइयाँ --- को छोटी, \omterm{prop:g4:journey-of-food:absorb}{रक्त} के लिए
तैयार इकाइयों में काटते हुए। दाँत और मथना बस टुकड़े छोटे करते हैं; रस
बदल देते हैं कि टुकड़े क्या हैं।
\end{solution}

\begin{solution}{exo:g7:digestion-and-nutrients:2}
शरीर की गरमाहट पर स्टार्च की लेई वाली दो नलियाँ: एक में ताज़ी \omterm{def:g4:journey-of-food:tract}{लार},
दूसरी में उबाली हुई \omterm{def:g4:journey-of-food:tract}{लार}। एक घंटे बाद ताज़ी वाली नली का स्टार्च शर्कराएँ
बन चुका है; उबाली वाली नली जस-की-तस है। निष्कर्ष: स्टार्च की तोड़-फोड़
अकेला रस करता है --- किसी मुँह की ज़रूरत नहीं --- और उबालना उसका काम
करने वाला पदार्थ नष्ट कर देता है।
\end{solution}

\begin{solution}{exo:g7:digestion-and-nutrients:3}
शर्कराएँ (\omterm{prop:g7:organs-at-work:bill}{ग्लूकोज़}) --- स्टार्च वाले और मीठे भोजनों से, जैसे रोटी;
शरीर बनाने वाले पदार्थ की इकाइयाँ --- माँस, मछली, \omterm{def:g2:animals-grow-up:born}{अंडों}, \omterm{def:g2:animals-grow-up:born}{दूध} की
चीज़ों, दालों से; चिकनाई की इकाइयाँ --- तेलों और मक्खन से; और साथ में
विटामिन, \omterm{prop:g4:what-plants-need:needs}{खनिज लवण} तथा पानी --- और वे बाक़ी चीज़ों के अलावा फलों और
सब्ज़ियों से।
\end{solution}

\begin{solution}{exo:g7:digestion-and-nutrients:4}
विटामिन, \omterm{prop:g4:what-plants-need:needs}{खनिज लवण} और पानी जैसे हैं वैसे ही पार हो जाते हैं। रेशा हमारे
रसों को पूरी तरह हरा देता है: बिना कटे वह कुछ भी पार नहीं करता और
बचे-खुचे के उपयोगी बुहारी वाले ढेर के रूप में आगे बढ़ जाता है।
\end{solution}

\begin{solution}{exo:g7:digestion-and-nutrients:5}
\omterm{def:g4:journey-of-food:tract}{छोटी आँत} की भीतरी दीवार का उँगली जैसा एक उभार, लगभग एक मिलीमीटर ऊँचा,
और लाखों में से एक। हर एक अपनी ही महीन रक्त-वाहिकाएँ लपेटे रहता है।
\end{solution}

\begin{solution}{exo:g7:digestion-and-nutrients:6}
\omterm{prop:g7:digestion-and-nutrients:absorption}{यकृत} तक --- यानी छाँटने-और-जमा करने वाले अंग तक। एक सेवा: खाने के
\omterm{prop:g7:organs-at-work:bill}{ग्लूकोज़} का एक हिस्सा बाद में छोड़ने के लिए जमा कर लेना।
\end{solution}

\begin{solution}{exo:g7:digestion-and-nutrients:7}
ढुलाई: भोजन के बड़े घटक आँत की दीवार पार नहीं कर सकते --- वही अलमारी
और झिरी-दरवाज़ा। और दोबारा जुड़ाई: शरीर अपना पदार्थ छोटी इकाइयों से
बनाता है, इसलिए उसे पुर्ज़े चाहिए, फ़र्नीचर नहीं --- यानी सफ़र के लिए
भी खुला बंडल, और दोबारा बनाने के लिए भी।
\end{solution}

\begin{solution}{exo:g7:digestion-and-nutrients:8}
वह एक निष्पक्ष परीक्षण का \omterm{met:g4:what-plants-need:fairtest}{नियंत्रण} है: रस के काम करने की हालत के सिवा
हर बात में एक जैसा। उसका बिना बदला स्टार्च तोड़-फोड़ का ज़िम्मा ख़ास
तौर पर ताज़ी \omterm{def:g4:journey-of-food:tract}{लार} पर डालता है --- और यह भी दिखाता है कि काम करने वाला
पदार्थ गरमी से नाज़ुक है, उबलने से किसी नाज़ुक मशीन की तरह नष्ट हो
जाता है।
\end{solution}

\begin{solution}{exo:g7:digestion-and-nutrients:9}
तीनों ही बड़ी, पतली, नम और \omterm{prop:g4:journey-of-food:absorb}{रक्त} से अस्तर की हुई सरहद-पार की सतहें हैं।
माल: फेफड़ों की कोठरियों और \omterm{def:g7:how-animals-breathe:four}{क्लोमों} के \omterm{def:g1:animals-around-us:covering}{पंखों} के लिए \omterm{def:g7:what-is-respiration:respiration}{ऑक्सीजन} तथा \omterm{def:g7:what-is-respiration:respiration}{कार्बन
डाइऑक्साइड}; और \omterm{prop:g7:digestion-and-nutrients:absorption}{रसांकुरों} के लिए \omterm{def:g7:digestion-and-nutrients:families}{पोषक तत्व}। एक विशेष-सूची, तीन सरहदें।
\end{solution}

\begin{solution}{exo:g7:digestion-and-nutrients:10}
\omterm{def:g2:animals-grow-up:born}{अंडा}: शरीर बनाने वाले पदार्थ का कुल। तोड़-फोड़ \omterm{def:g4:journey-of-food:tract}{आमाशय} के अम्लीय रस में
शुरू होती है (शीशे की आँत में उन टुकड़ों की नियति), और \omterm{def:g4:journey-of-food:tract}{छोटी आँत} में
पूरी होती है: शरीर बनाने वाली छोटी इकाइयाँ। सरहद-पार: \omterm{prop:g7:digestion-and-nutrients:absorption}{रसांकुर}, \omterm{prop:g4:journey-of-food:absorb}{रक्त},
\omterm{prop:g7:digestion-and-nutrients:absorption}{यकृत}। पहुँचाना: भरती हुई खरोंच तक, जहाँ बँटती हुई त्वचा-कोशिकाएँ उन
इकाइयों को दोबारा बनाने के सामान के तौर पर लेती हैं।
\end{solution}

\begin{solution}{exo:g7:digestion-and-nutrients:11}
शरीर बनाने वाले पदार्थ की इकाइयाँ कम पड़ जाती हैं। सबसे पहले वह सब
भुगतता है जो बन रहा है और नया हो रहा है: ख़ुद बढ़ना, \omterm{def:g3:bones-and-muscles:muscle}{पेशियों} की
देखभाल, मरम्मतें, और लगातार नई होती परतें तथा \omterm{prop:g4:journey-of-food:absorb}{रक्त} के घटक --- ईंधन
वाले \omterm{def:g7:digestion-and-nutrients:families}{पोषक तत्व} उनकी जगह नहीं ले सकते, क्योंकि ईंटें और जलावन आपस में
बदले नहीं जाते।
\end{solution}

\begin{solution}{exo:g7:digestion-and-nutrients:12}
मिसाल के लिए: खाना फ़र्नीचर की तरह भीतर आता है --- स्टार्च, शरीर बनाने
वाला पदार्थ, चिकनाई। ठिकाने-दर-ठिकाने उसके रस हर एक को उसकी इकाइयों में
काटते हैं: शर्कराएँ, शरीर बनाने वाली इकाइयाँ, चिकनाई की इकाइयाँ ---
यानी पहुँचाने की सूची। \omterm{prop:g7:digestion-and-nutrients:absorption}{रसांकुरों} पर वह सूची \omterm{prop:g4:journey-of-food:absorb}{रक्त} में उतर जाती है, \omterm{prop:g7:digestion-and-nutrients:absorption}{यकृत}
छाँटता, जमा करता और छोड़ता है, और \omterm{prop:g4:journey-of-food:absorb}{रक्त} बाँटता है: काम करती \omterm{def:g3:bones-and-muscles:muscle}{पेशियों} को
ईंधन, उसारी की जगहों को ईंटें, और बाक़ी भंडारों को। जो थाली पर भोजन बनकर
बैठा था, वह घंटों में पतों वाली सामग्री बनकर घूम रहा होता है।
\end{solution}

\begin{solution}{pb:g7:digestion-and-nutrients:1}
\textbf{1.} नली म: स्टार्च ग़ायब, शर्कराएँ हाज़िर। म$'$: स्टार्च
जस-का-तस। नली प: टुकड़ों पर हमला --- सिकुड़े हुए, मुलायम। प$'$: टुकड़े
जस-के-तस।

\textbf{2.} \omterm{def:g4:journey-of-food:tract}{लार} का काम करने वाला पदार्थ स्टार्च को शर्कराओं में तोड़
देता है। म$'$ \omterm{met:g4:what-plants-need:fairtest}{नियंत्रण} है: उबाले हुए रस के सिवा हर बात में एक जैसी होने
के नाते वह साबित करती है कि तोड़ने वाली ताज़ी \omterm{def:g4:journey-of-food:tract}{लार} ही थी, नली की और कोई
चीज़ नहीं।

\textbf{3.} उसने \omterm{def:g2:animals-grow-up:born}{अंडे} के शरीर बनाने वाले पदार्थ की तोड़-फोड़ शुरू कर दी
है --- उसे उसकी छोटी इकाइयों की ओर काटते हुए, यानी शरीर बनाने वाले पोषक
तत्वों की ओर। (साफ़ होते और सिकुड़ते टुकड़े वही कटाई हैं, \omterm{def:g1:the-five-senses:sense}{आँखों} से दिखी
हुई।)

\textbf{4.} यह कि असली काम रासायनिक है: किसी भी नली में न दाँत थे न
मथना, फिर भी तोड़-फोड़ वहीं-वहीं चली जहाँ सही रस सही भोजन से मिला।
यांत्रिक तोड़ना बस रसायन की सेवा करता है।

\textbf{5.} रस काम करने वाले पदार्थ हैं, और उनकी काम करने की
\omterm{def:g6:exploring-our-environment:conditions}{परिस्थितियाँ} हैं: ठंडे कर दो तो वे सुस्ता जाते हैं --- \omterm{def:g6:decomposers-and-soil:decomposers}{अपघटक} से लेकर
ख़मीर तक, हर मज़दूर की तरह। वह नली बिना पचे बैठी रहती, और सिर्फ़ यह
साबित करती कि बर्फ़ीले डिब्बे रसायन को ठहरा देते हैं।

\textbf{6.} यह कि हर रस की अपनी ख़ासियतें और अपने ही ठिकाने की
\omterm{def:g6:exploring-our-environment:conditions}{परिस्थितियाँ} हैं: एक ही नली में मिला दो तो कोई भी अपने सबसे अच्छे पर
काम नहीं करता, और नमूना नली-तंत्र का वह ठिकाना-दर-ठिकाना क्रम गँवा देता
है --- यानी ठीक वही चीज़, जो उसे सिखानी है।

\textbf{7.} अवशोषण करने वाली दीवार: यानी \omterm{prop:g7:digestion-and-nutrients:absorption}{रसांकुरों} के क़ालीन वाली एक
सतह। उसका सामान पतला होना चाहिए (ताकि पार करना आसान हो) और छोटी इकाइयों
के लिए चुनिंदा-पर-पारगम्य --- और शरीर में उसके पीछे रक्त-वाहिकाएँ भी,
जो पार करने वालों को आगे ले जाएँ।

\textbf{8.} चोकर जस-का-तस: हमारा कोई रस रेशा नहीं काटता। मतलब: नमूना
रेशे की असली नियति की ठीक भविष्यवाणी करता है --- कुछ भी पार न करना,
बुहारते हुए आगे बढ़ जाना --- यानी वह अपवाद, जो तंत्र की हदों का नक़्शा
बना देता है।

\textbf{9.} \omterm{prop:g7:digestion-and-nutrients:absorption}{रसांकुरों} की पतली दीवारों के पार; \omterm{prop:g4:journey-of-food:absorb}{रक्त} में; \omterm{prop:g7:digestion-and-nutrients:absorption}{यकृत} से होकर,
जो एक हिस्सा जमा कर लेता है; और वहाँ से आम बँटवारे पर शरीर की \omterm{def:g5:first-look-at-cells:cell}{कोशिकाओं}
तक --- जहाँ कोई काम करती \omterm{def:g3:bones-and-muscles:muscle}{पेशी} उन्हें उसी घंटे जला डालती है।

\textbf{10.} \omterm{def:g6:exploring-our-environment:conditions}{परिस्थितियाँ} ख़ुद मशीनरी का हिस्सा हैं: हर ठिकाना अपनी
तरह सधा है --- यहाँ अम्ल, वहाँ बेअसर किया हुआ --- इसलिए रसों को भोजन से
सही क्रम में मिलना ही चाहिए, और ठिकानों से होकर निगलना ठीक यही देता है।
नली-तंत्र एक बनाने की लाइन है, बाल्टी नहीं।

\textbf{11.} क्योंकि \omterm{def:g4:journey-of-food:digestion}{पाचन} के मज़दूर पदार्थ हैं, अंग नहीं: जैसा चौकी
दिखाती है, कोई रस अपने भोजन की तोड़-फोड़ वहीं कर देता है जहाँ दोनों सही
\omterm{def:g6:exploring-our-environment:conditions}{परिस्थितियों} में मिल जाएँ। इसलिए ग़ायब रस को सही ठिकाने तक पहुँचाने वाला
कैप्सूल उस ठिकाने का रसायन कर सकता है।

\textbf{12.} ``हर रस ने अपने ठिकाने पर भोजन को पोषक तत्वों में काटा ---
और नियंत्रण-नलियों ने साबित किया कि काटने का काम रस ही कर रहे थे, और
कोई नहीं।''
\end{solution}
